\subsection{摘要}\label{abstract}

AI 的进步并非由单一模型驱动，而在于能够持续改进当前模型的水平。要做到这一点，需要将模型开发视为一个系统级优化问题，其解决方案是构建一台用于快速改进的爬山机器（hill-climbing machine）。我们的流程包括一个以规模化为核心的预训练建模决策框架，以及一套稳健的强化学习配方与基础设施，能够维持长期、对数线性（log-linear）的性能提升。首个采用我们流程开发的模型是 MAI-Thinking-1，一个 35B 激活参数 / 1T 总参数的 MoE 模型，在 STEM 推理与编程任务上位居同类规模模型中的最强之列（例如，SWE-Bench Pro 上 52.8\%，AIME 2025 上 97.0\%，LiveCodeBench v6 上 87.7\%）。MAI-Thinking-1 完全从零训练，仅使用干净、企业级的数据，不依赖任何第三方模型的蒸馏。在本技术报告中，我们提供了对 MAI-Thinking-1 开发过程的深入剖析。通过分享我们的技术细节与经验教训，我们希望推动一种透明、以科学为导向的方法，促进 AI 的进一步发展。

\subsection{1 引言}\label{introduction}

我们介绍 MAI-Thinking-1，一个由微软 AI（MAI）从零开发的高性能推理模型，在 STEM 推理与编程任务上与同规模模型具有竞争力。我们的预训练聚焦于一种简单的规模化方法，强调以经验驱动的迭代改进来优化架构与数据；同时，我们的强化学习（RL）框架针对在数千步内持续的对数线性（log-linear）爬升进行了优化（图 1）。

我们遵循三条主要设计原则来开发模型：第一，能力应当被习得，而非被继承。虽然通过蒸馏模仿而来的智能获取速度更快，但它缺乏长期、持久爬升所必需的可引导性与稳健性。第二，简单性才是可持续的。我们偏好简单、可扩展的配方；干净、可信的数据；以及透明的、共同支撑从零爬升的基础设施。第三，科学严谨，不走捷径。每一项决策都必须能够通过数据驱动的梯级、消融实验与评估加以验证，从而揭示通往顶点的可靠路径。\footnote{Correspondence should be sent to mai-technical-report@microsoft.com. 引用格式见附录 A。}

\pandocbounded{\includegraphics[keepaspectratio,alt={image}]{images/2c9932a032103e844bd44827bf1431000b9f90c1767493aebceb3678ddea6736.jpg}}

\pandocbounded{\includegraphics[keepaspectratio,alt={image}]{images/5ca97040766ac1806564c80a7b670255903459a28f7bca2114edf15f10908c52.jpg}}

图 1. MAI-Thinking-1 在强化学习期间的性能。STEM 爬升（左与中）：我们以 STEM 为重点的爬升过程中的评估性能，评测使用 AIME 25 和 LiveCodeBench v6 的困难子集（2025 年 1 月之后）。对两条曲线，我们展示相邻 3 个检查点的 pass@1 平均值。智能体爬升（右）：在我们以代码为主的智能体爬升期间，SWE-bench Verified 上的 pass@1 性能。

\pandocbounded{\includegraphics[keepaspectratio,alt={image}]{images/c39dab32ae9864ffb503f1ed3b63e170fded6ec5d1bd6ac020a23361e9fe6ec6.jpg}}

MAI-Thinking-1 在 30T tokens 上完成预训练，这些 tokens 提取自公开可用与授权的人类生成数据的混合，涵盖网页数据、公开 GitHub 代码、书籍、学术论文、新闻、多语言文本和领域特定材料。每个数据源都在内部从始至终完成处理。我们选择在预训练中不使用任何由语言模型生成的合成数据，并努力避免与移除所收集数据源中的 AI 生成内容。在预训练中，我们不使用任何开源训练数据集，并从训练数据中去除常见的机器学习数据库（decontaminate）。在中期训练（mid-training）中，我们进一步强化 STEM、数学和编程能力，为推理 RL 爬升奠定坚实基础。中期训练之后，MAI-Thinking-1 的最大上下文长度达到 256K。预训练和中期训练赋予我们的基础模型广泛的预测能力与知识，但它们并未规定模型应如何表现、如何解决长程任务，或如何分配推理时的计算资源。

在强化学习（RL）爬升阶段，我们教会模型推理与回应。模型学会利用思维链（CoT）应对任务特定的反馈、使用外部工具与环境交互，以及遵循人类偏好与安全信号。我们的 RL 爬升从零开始，在没有任何先验推理轨迹接触的情况下学习推理。稳健的 RL 配方、自蒸馏（self distillation）与基础设施改进，使我们能够将 RL 运行持续数千步。利用这一配方，我们训练了三个领域专用的专家模型：一个用于 STEM 推理，一个用于智能体编程与工具使用，一个用于有用性与安全性。

我们将 MAI-Thinking-1 的性能与一系列开放权重与封闭权重的前沿模型，在公开基准和人类并排对比评估中进行比较。这些评估涵盖广泛的领域，凸显了模型在不同领域中的通用性：STEM、智能体编程、知识、指令遵循、长上下文、安全、健康与工具调用。MAI-Thinking-1 在 SWE-Bench Pro 上取得 52.8\%，在 AIME 2025 上取得 97.0\%，在 AIME 2026 上取得 94.5\%，在 LiveCodeBench v6 上取得 87.7\%，并在广泛的基准上与 Sonnet 4.6 具有竞争力。

开发 MAI-Thinking-1 时的一个关键技术挑战，是在安全性（要求模型拒绝用户请求）与有用性（要求模型遵从用户）之间取得平衡。一个安全的模型需要提供有帮助的回应，同时符合我们的安全政策与标准。我们开发内部安全基准来锚定进展，并将有用性与安全性训练纳入 RL 爬升。作为负责任部署的一部分，我们在整个模型开发过程中持续进行红队测试（red-team），在发布前发现并修复漏洞。

MAI-Thinking-1 是首个使用我们的爬山机器开发的模型：它整合了数据管线构建、训练基础设施、强化学习环境与奖励、评估套件与安全测试，将模型开发转化为针对特定领域的经验优化循环。爬山机器使我们在推进 AI 的同时，将进展自下而上地锚定于人类需求之上。
